\documentclass[12pt, twoside, openany]{book}

\usepackage{mathptmx} % Contiene una fuente similar a Times New Roman

\usepackage[spanish, es-tabla]{babel} % Permite escritura en castellano
\usepackage[utf8]{inputenc} % Permite utilizar caracteres UTF8

\usepackage{graphicx} % Para la inclusión de gráficos e imágenes
\graphicspath{ {images/} } % Ruta para buscar las imágenes
\usepackage[a4paper,top=30mm,left=30mm,right=25mm,bottom=25mm,headheight=20mm]{geometry} % Configuración de los margenes de la página

% Paquetes para que funcione el formato.
\usepackage{titlesec}
\usepackage{setspace}
\usepackage{ragged2e}
\usepackage{fancyhdr}
\usepackage{lastpage}
\usepackage{stackengine}
\usepackage{array}
\usepackage{url}
\usepackage{float}

\usepackage{hyperref} % Paquete para que las referencias funcionen, y permite introducir links
\usepackage{xcolor} % Paquete para trabajar con colores (fondo de celdas, color del texto...)
\usepackage{pdfpages} % Paquete para incluir páginas de otro PDF directamente
\usepackage{lipsum} % Paquete para generación de texto
% Se define un color gris desde su código RGB
\definecolor{gris}{RGB}{220,220,220}

\setcounter{secnumdepth}{3} % Para permitir numerar las sub-subsecciones

% Modifica el nombre de los índices al castellano
\addto\captionsspanish{
  \renewcommand{\contentsname}{Índice de contenido}
  \renewcommand{\listfigurename}{Índice de figuras}
  \renewcommand{\listtablename}{Índice de tablas}
}

% Formateo de los nombres de los apartados:
\titleformat{\chapter}[block]
  {\normalfont\Huge\bfseries\singlespacing}{\thechapter.}{1em}{\Huge}
\titlespacing*{\chapter}{0pt}{-62pt}{0pt}

\titleformat{\section}[block]
  {\normalfont\large\bfseries}{\thesection.-}{4pt}{\large}
\titlespacing*{\section}{0pt}{\baselineskip}{0pt}

\titleformat{\subsection}[block]
  {\normalfont\normalsize\bfseries}{\thesubsection.-}{4pt}{\normalsize}
\titlespacing*{\subsection}{0pt}{0pt}{0pt}

\titleformat{\subsubsection}[block]
  {\normalfont\normalsize\bfseries}{\thesubsubsection.-}{4pt}{\normalsize}
\titlespacing*{\subsubsection}{0pt}{0pt}{0pt}

\def\tablename{Tabla}

% Formato de header y footer
\fancyhf{}
\fancyhead[L]{\includegraphics[height=16mm]{style/logo_epi.png}
  \hspace{1em} \Longstack[l]{
    \textbf{UNIVERSIDAD DE OVIEDO} \newline
    \textbf{Escuela Politécnica de Ingeniería de Gijón}}}
\fancyhead[R]{\bfseries{Hoja \thepage \hspace{1pt} de \pageref{LastPage}}}
\fancyfoot[C]{Nombre Apellido1 Apellido2}
\renewcommand{\headrulewidth}{0pt}
\renewcommand{\footrulewidth}{0.4pt}

% Reescritura del plain para que salgan las mismas cabeceras en todas las paginas
\fancypagestyle{plain}{%
  \fancyhf{}
  \fancyhead[L]{\includegraphics[height=16mm]{style/logo_epi.png}
    \hspace{1em} \Longstack[l]{
      \textbf{UNIVERSIDAD DE OVIEDO} \newline
      \textbf{Escuela Politécnica de Ingeniería de Gijón}}}
  \fancyhead[R]{\bfseries{Hoja \thepage \hspace{1pt} de \pageref{LastPage}}}
  \fancyfoot[C]{Nombre Apellido1 Apellido2}
  \renewcommand{\headrulewidth}{0pt}% default is 0pt
  \renewcommand{\footrulewidth}{0.4pt}% default is 0pt
}

\pagestyle{fancy}

\restylefloat{table}




\begin{document}

\rmfamily % Fuente tipo Romana

% Portada de la memoria
\begin{titlepage}
    \centering
    \bfseries {
        \null
        \vspace{0cm}
        \begin{table}[h]
            \centering
            \begin{tabular}{m{10cm} m{1cm} m{3cm}}
                \vspace{0.2cm}
                \includegraphics[width=86mm]{style/logo_uniovi.jpg} &  & \vspace{1.52mm} \includegraphics[width=23mm]{style/logo_epi.png} \\
            \end{tabular}
        \end{table}

        \vspace{3\baselineskip}

        \Large{
            ESCUELA POLITÉCNICA DE INGENIERÍA DE GIJÓN \\ \vspace{3\baselineskip}
        }
        \large{
            GRADO EN INGENIERÍA INFORMÁTICA EN TECNOLOGÍAS DE LA INFORMACIÓN \\ \vspace{3\baselineskip}
            TÍTULO WIP \\
            \vspace{3\baselineskip}
            Mier Montoto, Juan Francisco \\
            UO283319 \\
            \vspace{2\baselineskip}
            FECHA:\@ diciembre 2023
        }
    }
\end{titlepage}


% Índice de contenido
\addcontentsline{toc}{chapter}{Índice de contenido} % Añade la referencia al índice de contenido
\tableofcontents
\newpage

% Índice de figuras
\addcontentsline{toc}{chapter}{Índice de figuras}  % Añade la referencia al índice de contenido
\listoffigures
\newpage

% Índice de tablas
\addcontentsline{toc}{chapter}{Índice de tablas} % Añade la referencia al índice de contenido
\listoftables


\justify % Texto justificado
\setlength{\parskip}{\baselineskip} % Separación entre párrafos de 1 linea
\onehalfspacing % Interlineado de 1,5

%% Incluir aquí el contenido de la memoria
\chapter{Introducción}


\chapter{Conclusiones}


%% Esto incluirá la bibliografía correctamente en nuestro trabajo
\newpage % En una nueva página
\addcontentsline{toc}{chapter}{Bibliografía} % Añade la referencia al índice de contenido

\bibliographystyle{ieeetr} % Define el estilo de la bibliografía
\bibliography{Bibliografia} % Indica el archivo que contiene la colección de citas

\end{document}
